\sloppy 
\phantomsection

\addcontentsline{toc}{chapter}{Index} 


\addcontentsline{toc}{section}{\lsNameIndexTitle}
\ohead{\lsNameIndexTitle}
\printindex 
\cleardoublepage


\phantomsection
\addcontentsline{toc}{section}{Index of books}
\ohead{Index of books}
\printindex[ban]
\ohead{}

\section*{Abbreviations used in book index}

\lehead{Book index}
\rohead{Book index}
\begin{description}

  \setlength{\itemsep}{0em}
%   \setlength{\parsep}{0pt}
%   \setlength{\topsep}{0pt}
\item[\bfseries\upshape AbB 1]Kraus, Fritz R. 1964. \emph{Briefe aus dem British Museum (CT 43 \& 44)} (Altbabylonische Briefe in Umschrift und Übersetzung 1). Leiden: Brill.
\item[\bfseries\upshape AbB 7]Kraus, Fritz R. 1977. \emph{Briefe aus dem British Museum (CT 52)} (Altbabylonische Briefe in Umschrift und Übersetzung 7). Leiden: Brill.
\item[\bfseries\upshape ABL 1]Harper, Robert F. 1892. \emph{Assyrian and Babylonian letters} (Assyrian and Babylonian Letters 1). London \& Chicago: Chicago University Press.
\item[\bfseries\upshape AEW]de Vries, Jan. 1977. \emph{Altnordisches Etymologisches Wörterbuch.} 2nd ed. Leiden: Brill.
\item[\bfseries\upshape ALEW ]Hock, Wolfgang, Rainer Fecht, Anna H. Feulner, Eugen Hill and Dagmar S. Wodtko. N.d. \emph{Altlitauisches etymologisches Wörterbuch (ALEW).} Version 2.0.  https://alew.uni-frankfurt.de.
\item[\bfseries\upshape AnGr]Noreen, Adolf. 1923. \emph{Altnordische Grammatik I.} 2nd ed. Halle: Niemeyer.
\item[\bfseries\upshape ARM 1]Dossin, Georges. 1950. \emph{Correspondence de Šamši-Addu} (Archives Royales de Mari 1). Paris: Imprimerie nationale.
\item[\bfseries\upshape BT]Bosworth, Joseph \& Thomas Toller. 1996. \emph{An Anglo-Saxon Dictionary.} London: Clarendon.
\item[\bfseries\upshape CDA]Black, Jeremy A., Andrew R. George \& John Nicholas Postgate. 2000. \emph{A concise dictionary of Akkadian.} 2nd revised ed., vol. 5 (Santag: Arbeiten und Untersuchungen zur Keilschriftkunde). Wiesbaden: Harrassowitz Verlag.
\item[\bfseries\upshape CHD]Goedegebuure, Petra M., Hans G. Güterbock, Harry A. Hoffner \& Theo P. J. van den Hout. 1980–. \emph{The Hittite dictionary of the Institute for the Study of Ancient Cultures of the University of Chicago} (Chicago Hittite dictionary). Chicago: The Oriental Institute of the University of Chicago.
\item[\bfseries\upshape CT 28]Wallis Budge, Ernest A. 1910. \emph{Cuneiform texts from Babylonian tablets in the British Museum}, vol. 28. London: Harrison \& sons.
\item[\bfseries\upshape CT 34]Wallis Budge, Ernest A. 1914. \emph{Cuneiform texts from Babylonian tablets in the British Museum}, vol. 34. London: Harrison \& sons.
\item[\bfseries\upshape CT 39]Gadd, Cyril J. 1926. \emph{Cuneiform texts from Babylonian tablets in the British Museum}, vol. 39. London: Harrison \& sons.
\item[\bfseries\upshape CV]Cleasby, Richard \& Guðbrandur Vigfússon. 1874. \emph{An Icelandic-English dictionary.} Oxford: Clarendon.
\item[\bfseries\upshape EDG]Beekes, Robert. 2010. \emph{Etymological dictionary of Greek.} Leiden: Brill.
\item[\bfseries\upshape eDiAna]Hackstein, Olav, Jared Miller, Elisabeth Rieken \& Ilya Yakubovich (eds.). N.d. \emph{Digital philological-etymological dictionary of the minor Anatolian corpus languages.} https://www.ediana.gwi.uni-muenchen.de (05.09.2025).
\item[\bfseries\upshape eDIL]N. A. 2019. \emph{An electronic dictionary of the Irish Language, based on the Contributions to a Dictionary of the Irish Language.} https://dil.ie/ (23 November, 2023).
\item[\bfseries\upshape EDL]de Vaan, Michiel. 2008. \emph{Etymological dictionary of Latin and the other Italic languages.} Leiden: Brill.
\item[\bfseries\upshape EDPG]Kroonen, Guus. 2013. \emph{Etymological dictionary of Proto-Germanic.} Leiden: Brill.
\item[\bfseries\upshape ESJS]Havlová, Eva, Ilona Janyšková and Adolf Erhart (ed.). 1989–2023. \emph{Etymologický slovník jazyka staroslověnského.} Praha: Academia - Nakladatelství Československé Akademie věd.
\item[\bfseries\upshape EVP]Renou, Louis. 1955–1969. \emph{Études védiques et pāṇinéennes}, vols. I–XVII. Paris: Institut de Civilisation Indienne.
\item[\bfseries\upshape EWA]Lühr, Rosemarie, Albert L. Llyod \& Otto Springer (eds.). 1988. \emph{Etymologisches Wörterbuch des Althochdeutschen.} Göttingen: Vandenhoeck \& Ruprecht.
\item[\bfseries\upshape EWAia I]Mayrhofer, Manfred. 1992. \emph{Etymologisches Wörterbuch des Altindoarischen}, vol. I. Heidelberg: Winter.
\item[\bfseries\upshape EWAia II]Mayrhofer, Manfred. 1996. \emph{Etymologisches Wörterbuch des Altindoarischen}, vol. II. Heidelberg: Winter.
\item[\bfseries\upshape EWGP]Heidermanns, Frank. 1993. \emph{Etymologisches Wörterbuch der germanischen Primäradjektive.} Berlin: De Gruyter.
\item[\bfseries\upshape GDLI]Battaglia, Salvatore (ed.). 1961–2009. \emph{Grande dizionario della lingua italiana.} Turin: Utet. http://www.gdli.it.
\item[\bfseries\upshape GPC]Thomas, Richard J., Gareth A. Bevan \& Patrick J. Donovan (eds.). 1967–2002. \emph{Geiriadur Prifysgol Cymru, A Dictionary of the Welsh Language.} 4 vols. Caerdydd: Gwasg Prifysgol Cymru. http://welsh-dictionary.ac.uk/gpc/gpc.html   (23 November, 2023).
\item[\bfseries\upshape GRADIT]De Mauro, Tullio (ed.). 1999–2007. \emph{Grande dizionario italiano dell’uso.} Turin: Utet.
\item[\bfseries\upshape HW]Friedrich, Johannes \& Annelies Kammenhuber. 1975–1984. \emph{Hethitisches Wörterbuch.} 2. völlig neu bearbeitete Auflage auf der Grundlage der edierten hethitischen Texte. Heidelberg: Winter.
\item[\bfseries\upshape IEW]Pokorny, Julius. 1959–1969. \emph{Indogermanisches etymologisches Wörterbuch.} Bern: Francke.
\item[\bfseries\upshape KAH 2]Schroeder, Otto. 1922. \emph{Keilschrifttexte historischen Inhalts (KAH)}, vol. 2 (Wissenschaftliche Veröffentlichungen der deutschen Orient-Gesellschaft 37). Leipzig: J. C. Hinrichs’sche Buchhandlung.
\item[\bfseries\upshape KAR 1]Ebeling, Erich. 1919. \emph{Keilschrifttexte aus Assur religiösen Inhalts}, vol. 1 (Wissenschaftliche Veröffentlichungen der deutschen Orient-Gesellschaft 28). Leipzig: J. C. Hinrichs’sche Buchhandlung.
\item[\bfseries\upshape KAR 2]Ebeling, Erich. 1920. \emph{Keilschrifttexte aus Assur religiösen Inhalts}, vol. 2 (Wissenschaftliche Veröffentlichungen der deutschen Orient-Gesellschaft 34). Leipzig: J. C. Hinrichs’sche Buchhandlung.
\item[\bfseries\upshape KEWA III]Mayrhofer, Manfred. 1976. \emph{Kurzgefasstes etymologisches Wörterbuch des Altindischen}, vol. III. Heidelberg: Universitätsverlag Winter.
\item[\bfseries\upshape Labat TDP]Labat, René. 1951. \emph{Traité akkadien de diagnostics et prognostics médicaux}, vol. 1 (Collections de travaux de l’académie internationale d’histoire des sciences 7). Leiden: Brill.
\item[\bfseries\upshape Lambert BWL]Lambert, Wilfred G. 1996. \emph{Babylonian wisdom literature.} Winona Lake: Eisenbrauns.
\item[\bfseries\upshape LEIA]Vendryes, Joseph, E. Bachellery \& Pierre-Yves Lambert (eds.). 1959-1996. \emph{Lexique etymologique de l’irlandais ancien.} 7 vols. Dublin: Institute for Advanced Studies.
\item[\bfseries\upshape LIV2]Rix, Helmut. 2001. \emph{Lexikon der indogermanischen Verben.} Bearbeitet von Martin Kümmel, Thomas Zehnder, Reiner Lipp und Brigitte Schirmer. 2nd ed. Wiesbaden: Reichert.
\item[\bfseries\upshape LIV2 Addenda]Kümmel, Martin J. 2023. \emph{Addenda und Corrigenda zu LIV2.} Version 11.08.2023 16:56. http://www.martinkuemmel.de/liv2add.pdf.
\item[\bfseries\upshape LKG]Ulvydas, Kazys (ed.). 1971. \emph{Lietuvių kalbos gramatika.} 2: Morfologija. Vilnius: Mintis.
\item[\bfseries\upshape LKŽ]Lietuvių Kalbos Institutas. 1941–2002. \emph{Lietuvių kalbos žodynas.} Vilnius: Lietuvių kalbos ir literatūros institutas. https://www.lkz.lt.
\item[\bfseries\upshape LP]Wills, Tarrin (ed.). 2024. \emph{Lexicon poeticum.} https://lexiconpoeticum.org  (10 October, 2024).
\item[\bfseries\upshape MhdGr]Klein, Thomas, Hans-Joachim Soms \& Klaus-Peter Wegera. 2009. \emph{Mittelhochdeutsche Grammatik.} Teil 3: Wortbildungslehre. Boston: De Gruyter.
\item[\bfseries\upshape MLLVG]Sokols, Ēvalds, Anna Bergmane, Rūdolfs Grabis \& Milda Lepika (eds.). 1959. \emph{Mūsdienu latviešu literārās valodas gramatika} 1: Fonētika un morfoloģija. Rīga: Latvijas PSR Zinātņu akadēmijas Latviešu valodas institūts.
\item[\bfseries\upshape NIL]Wodtko, Dagmar S., Britta Irslinger \& Carolin Schneider (eds.). 2008. \emph{Nomina im indogermanischen Lexikon.} Heidelberg: Winter.
\item[\bfseries\upshape OgnS]Fritzner, Johan. 1867. \emph{Ordbog over det gamle norske Sprog.} Kristiania: Feilberg \& Landmarks.
\item[\bfseries\upshape ONP]Sigurðardóttir, Aldís (ed.) 2005–. \emph{Old Norse prose dictionary.} https://onp.ku.dk/onp/onp.php.
\item[\bfseries\upshape RIMA 2]Grayson, A. Kirk. 1991. \emph{Assyrian Rulers of the Early First Millennium BC}, vol. 1, (1114–859 BC) (The Royal Inscriptions of Mesopotamia 2). Toronto: University of Toronto Press.
\item[\bfseries\upshape TIM 2]Cagni, Luigi. 1980. \emph{Briefe aus dem Iraq Museum} (Texts in the Iraq Museum (TIM) 2). Leiden: Brill.
\item[\bfseries\upshape WAP]Baetke, Walter. 1965–1968. \emph{Wörterbuch zur altnordischen Prosaliteratur.} Berlin: Akademie Verlag.
\item[\bfseries\upshape YOS 2]Lutz, Henry F. 1917. \emph{Early Babylonian letters from Larsa} (Yale Oriental Series 2). New Haven: Yale University Press.
\end{description}
\cleardoublepage

%
% \phantomsection
% \addcontentsline{toc}{section}{\lsLanguageIndexTitle}
% \ohead{\lsLanguageIndexTitle}
%\printindex[lan]
% \cleardoublepage
  
\phantomsection
\addcontentsline{toc}{section}{\lsSubjectIndexTitle}
\ohead{\lsSubjectIndexTitle}
\printindex[sbj]
\ohead{}

\phantomsection
\addcontentsline{toc}{section}{Word forms}
\ohead{Form index}
\printindex[wan]
\ohead{}

% \phantomsection
% \addcontentsline{toc}{section}{Word forms, sorted}
% \ohead{Form index 2}
% \printindex[van]



